\section{Experimental Setup}
\label{sec:setup}

\paragraph{Training.} Every learner trains on \texttt{random\_day} (24\,h, 288
decisions, randomized bursts, failures and flash crowds per episode) for
400{,}000 transitions from 16 vectorized environments, with checkpoints every
50{,}000 transitions. The two learners of a training root share the training
episode-seed ledger byte for byte, so roots are paired.

\paragraph{Evaluation protocol.} Seven scripted scenarios
(Table~\ref{tab:scenarios}) are evaluated with three disjoint seed sets that are
fixed in code before any new result was produced: \emph{validation} seeds
(2001--2005, 35 episodes) select a checkpoint (highest mean return; ties to the
earlier); \emph{test} seeds (3001--3020, 140 episodes) are used only for
reporting, for the selected and for the final checkpoint; and
\emph{diagnostic} seeds (4001--4003) are used only for structural diagnostics.
The closed study's selection seeds (101--105) and final-holdout seeds
(1001--1005) are used only to reproduce its numbers
(Section~\ref{sec:repro}). Evaluation is deterministic (masked argmax). Because
traffic is exogenous, every policy faces byte-identical traffic on a given
(scenario, seed), so all comparisons are paired.

\paragraph{Metric.} The primary metric is the undiscounted operational return
$R=\sum_t r_t$ of an episode, averaged over the 7 scenarios $\times$ 20 seeds.
Scenarios have 48--288 decisions, so the average weights long scenarios more;
we therefore also report per-scenario results and the operational components
(delivered ratio, SLA-violating demand-intervals, maximum utilization,
reroutes per hour, moved bandwidth).

\paragraph{Statistics.} We separate two sources of uncertainty. For a
\emph{fixed} policy, 95\,\% intervals are scenario-stratified bootstrap
intervals over test episodes (10{,}000 resamples). For a \emph{learner}, the
experimental unit is the training root: intervals are two-stage bootstrap
intervals (resample roots, then scenario-stratified episodes within roots), and
we also report a $t$-interval over root-level means and the number of roots on
which a difference has a given sign. Paired differences are reported in reward
units together with the paired standardized effect $d_z$. We do not use
significance tests as decision rules~\citep{agarwal2021precipice,patterson2024empirical}.
The two-stage bootstrap conditions on the roots that were drawn: with two or
three roots it stays close to the range of the root means and therefore
understates between-root uncertainty. For the five-root primary comparison we
report both intervals and lead with the $t$-interval; for ablations with two or
three roots (horizon sweep, delay, tuning) we report per-root values and sign
counts, and the bootstrap interval only as a description of episode-level
variability. Per-scenario intervals resample roots, then episodes within roots.

\paragraph{Hardware.} All new results were produced on a 4-core Intel Xeon
(2.1\,GHz) virtual machine without GPU, one thread per run (PyTorch 2.14 CPU,
SB3/SB3-contrib 2.9.0). The closed study trained on an RTX~4070 laptop GPU; its
checkpoints were not archived, so we retrained them (Section~\ref{sec:repro}).
